\section*{Chapter \ref{ch:b2:isotope-geochemistry} --- Isotope Geochemistry}

\begin{solution}{exo:b2:isotope-geochemistry:1}
Tristan: $(0.51249/0.512638 - 1)\times10^4 = -2.9$. Mangaia: $(0.51285/0.512638 -
1)\times10^4 = +4.1$. MORB: $+8.5$. The sources of Mangaia and of MORB, with positive
$\varepsilon_{\mathrm{Nd}}$, have had a higher Sm/Nd than \omterm{def:b2:planet-formation-meteorites:chondrule}{chondrites} on average; the
source of Tristan a lower one.
\end{solution}

\begin{solution}{exo:b2:isotope-geochemistry:2}
$\delta^{18}\mathrm{O}_{\mathrm{VPDB}} = (\num{2.0672e-3}/\num{2.0052e-3} - 1)\times1000 =
+30.9$ per mil on the \omterm{def:b2:isotope-geochemistry:standards}{VSMOW} scale. The calcite has $^{18}\mathrm{O}/^{16}\mathrm{O} =
\num{2.0672e-3}\times(1 - 0.0020) = \num{2.0631e-3}$, hence $\delta^{18}\mathrm{O} =
(\num{2.0631e-3}/\num{2.0052e-3} - 1)\times1000 = +28.9$ per mil (\omterm{def:b2:isotope-geochemistry:standards}{VSMOW}).
\end{solution}

\begin{solution}{exo:b2:isotope-geochemistry:3}
At \qty{1.0}{Ga}: $e^{\lambda t} - 1 = 0.006545$, and $0.512638 - 0.1966\times0.006545 =
0.511351$. At \qty{4.567}{Ga}: $e^{\lambda t}-1 = 0.030243$, and $0.512638 -
0.1966\times0.030243 = 0.506692$.
\end{solution}

\begin{solution}{exo:b2:isotope-geochemistry:4}
$\delta\mathrm{D} = 8\times(-18.0) + 10 = -134$ per mil. For the lake, $d = -42 -
8\times(-4.0) = -10$ per mil, below the line's 10: the water lies below the \omterm{prop:b2:isotope-geochemistry:mwl}{meteoric
water line}, as evaporated surface waters do (the light isotopes left first).
\end{solution}

\begin{solution}{exo:b2:isotope-geochemistry:5}
$Q = 10^4\times\num{6.524e-12}\times0.1966/0.512638 = \num{2.502e-8}$ per year, or
\qty{25.0}{\per\giga\year}. The granite has $\mathrm{Sm}/\mathrm{Nd} = 6.0/30 = 0.200$,
against $0.145/0.457 = 0.317$ for CI: $f = 0.200/0.317 - 1 = -0.37$. Over
\qty{0.5}{Gyr}: $\Delta\varepsilon = 25.0\times(-0.37)\times0.5 = -4.6$ units.
\end{solution}

\begin{solution}{exo:b2:isotope-geochemistry:6}
$\varepsilon \approx 3.5 + 25.0\times(-0.45)\times2.7 = 3.5 - 30.4 = -26.9$, against
$-27.3$: the approximation is 0.4 unit too high (not negative enough). Most of the
difference, about 0.3 unit, comes from $e^{\lambda t} > 1$: the parent, hence the
deficit in growth relative to CHUR, was more abundant in the past than today
(\cref{exo:b2:isotope-geochemistry:10}); the rest comes from rounding $f$ and the initial value.
\end{solution}

\begin{solution}{exo:b2:isotope-geochemistry:7}
With a mass fraction $x$ of granite, $\mathrm{Sr} = 320x + 130(1-x)$ and
$R = (320x\times0.7200 + 130(1-x)\times0.7030)/\mathrm{Sr}$. For $x = 0.05$: Sr
\qty{139.5}{ppm}, $R = 0.70495$; $x = 0.10$: 149.0, 0.70665; $x = 0.20$: 168.0,
0.70948. From the proof of the mixing proposition, $R\,\mathrm{Sr}$ is linear in Sr:
$R = 0.731632 - 3.7221/\mathrm{Sr}$, a straight line of slope $-3.722$ (ppm) and
intercept 0.7316 against $1/\mathrm{Sr}$. Check at $x = 0.05$: $0.731632 - 3.7221/139.5
= 0.70495$.
\end{solution}

\begin{solution}{exo:b2:isotope-geochemistry:8}
At \qty{20}{\degreeCelsius}: $\delta^{18}\mathrm{O} = 1000/1.0098 - 1000 = -9.70$ per
mil, $\delta\mathrm{D} = 1000/1.084 - 1000 = -77.5$ per mil, slope $77.5/9.70 = 7.98$.
At \qty{0}{\degreeCelsius}: $-11.56$ and $-99.9$ per mil, slope 8.64. Both are close to
the slope 8 of the \omterm{prop:b2:isotope-geochemistry:mwl}{meteoric water line}: the equilibrium fractionation of water
between liquid and vapour alone sets the ratio of the deuterium and oxygen-18
shifts near 8.
\end{solution}

\begin{solution}{exo:b2:isotope-geochemistry:9}
The crust holds $0.0051\times27/1.25 = 0.110$, 11 per cent of the neodymium. Mass
balance: $0.110\times(-15) + 0.890\,\varepsilon_m = 0$, so $\varepsilon_m = 1.65/0.890 =
+1.9$. With a depleted reservoir holding a share $x$ at $+8.5$ and the rest at 0:
$0.110\times(-15) + 8.5x = 0$, $x = 1.65/8.5 = 0.19$: about a fifth of the
neodymium.
\end{solution}

\begin{solution}{exo:b2:isotope-geochemistry:10}
The growth rate is $\lambda p\,e^{\lambda t}$, with $e^{\lambda t} \approx 1 + \lambda t$;
averaged over the interval from the separation $T$ to the present, the factor is
$1 + \lambda T/2$, so the neglected terms are of relative size $\lambda T/2$. Over the
age of the Earth, $\num{6.524e-12}\times\num{4.567e9}/2 = 0.015$, 1.5 per cent. For the
crust of \cref{exo:b2:isotope-geochemistry:6}, $\lambda T/2 = 0.0088$, and the error is $0.0088\times30.4 = 0.27$
unit. The slow change of the CHUR ratio by which $\varepsilon$ is divided adds
another, smaller term, of order $(p/R)\lambda T \times \varepsilon_0$, here 0.02 unit.
\end{solution}

\begin{solution}{exo:b2:isotope-geochemistry:11}
$K = 13.3/5.0 = 2.66$. With $s = 0.5$: $s' = 2.66\times0.5/(1 + 1.66\times0.5) =
1.33/1.83 = 0.73$. The difference $s' - s$ is largest where
$\mathrm{d}s'/\mathrm{d}s = K/(1+(K-1)s)^2 = 1$, at $s = (\sqrt K - 1)/(K-1) =
0.631/1.66 = 0.38$, where $s' = 0.62$: 0.24 apart.
\end{solution}

\begin{solution}{exo:b2:isotope-geochemistry:12}
Generation: $2.3$ to $2.6\times\qty{7.2e9}{km^3}/\qty{4.5e9}{yr} = 3.7$ to
\qty{4.2}{km^3} a year. Destruction: $1.3$ to $1.6\times7.2/4.5 = 2.1$ to
\qty{2.6}{km^3} a year. The present rates differ: additions (2.7) are below the model's
mean generation, removal (3.2) above its mean destruction. In the model the crust grew
early, when generation was fastest, and the present is a time of balance or slight
loss, as the subduction-zone budget says.
\end{solution}

\begin{solution}{pb:b2:isotope-geochemistry:1}
\textbf{1.} Mantle: $(0.51307/0.512638 - 1)\times10^4 = +8.4$. Sediment:
$(0.51141/0.512638 - 1)\times10^4 = -24.0$.

\textbf{2.} Kora $+2.8$, Vel $-4.3$, Asha $+4.7$.

\textbf{3.} With $^{87}\mathrm{Sr}/^{86}\mathrm{Sr}$ across and $\varepsilon_{\mathrm{Nd}}$ up:
mantle (0.70282, $+8.4$), Kora (0.70410, $+2.8$), Vel (0.70618, $-4.3$), sediment
(0.72000, $-24.0$), Asha (0.70290, $+4.7$). Only Vel, with a negative
$\varepsilon_{\mathrm{Nd}}$, lies on the enriched side of CHUR.

\textbf{4.} Below the array: at the strontium ratio of the \omterm{def:b2:isotope-geochemistry:reservoirs}{depleted mantle}, the array
would have $\varepsilon_{\mathrm{Nd}}$ near $+8$; Asha has $+4.7$. It resembles the
HIMU-type St.~Helena group (Mangaia, Tubuai, St.~Helena).

\textbf{5.} Melting changes the parent--daughter ratios but not, measurably, the
isotope ratios of strontium and neodymium (proposition on melting); the lavas are
young, so their ratios have not grown since eruption (question 6).

\textbf{6.} $\Delta R = 0.12\,(e^{\lambda\cdot10^6} - 1) = 0.12\times\num{6.52e-6} =
\num{7.8e-7}$, or $\num{7.8e-7}/0.512638\times10^4 = 0.015$ epsilon unit: far below the
differences between the islands, and below the precision of a measurement.

\textbf{7.} $\mathrm{Nd} = 40f + 0.75(1-f)$, and
$R = \bigl(40f\times0.51141 + 0.75(1-f)\times0.51307\bigr)/\mathrm{Nd}$.

\textbf{8.} $f = 0.5$ per cent: Nd \qty{0.9463}{ppm}, $\varepsilon = +1.6$; 1 per cent:
1.1425, $-2.9$; 2 per cent: 1.535, $-8.4$; 5 per cent: 2.7125, $-15.4$.

\textbf{9.} Sr 10.95, 11.90, 13.80 and \qty{19.50}{ppm}; $^{87}\mathrm{Sr}/^{86}\mathrm{Sr} =
0.70439$, 0.70571, 0.70780 and 0.71163.

\textbf{10.} $R\,\mathrm{Sr} = 10\times0.70282 + f(200\times0.72000 - 10\times0.70282)$ and
$\mathrm{Sr} = 10 + 190f$; eliminating $f$, $R\,\mathrm{Sr} = 7.0282 +
(\mathrm{Sr}-10)\times136.9718/190$, so $R = 0.720904 - 0.18084/\mathrm{Sr}$: a straight
line of slope $-0.1808$ (ppm) against $1/\mathrm{Sr}$. Check at $f = 1$ per cent:
$0.720904 - 0.18084/11.90 = 0.70571$.

\textbf{11.} $K = (10/0.75)/(200/40) = 13.33/5.0 = 2.67$. The mantle is relatively
richer in strontium, the sediment in neodymium, so the Nd ratio moves towards the
sediment faster: a little sediment shifts $\varepsilon_{\mathrm{Nd}}$ a long way before
$^{87}\mathrm{Sr}/^{86}\mathrm{Sr}$ moves much.

\textbf{12.} 53 times richer in Nd ($40/0.75$), 20 times in Sr. At $f = 1$ per cent the
sediment supplies $0.40/1.1425 = 35$ per cent of the neodymium, which moves
$\varepsilon_{\mathrm{Nd}}$ by $0.35\times(8.4 + 24.0) = 11.3$ units, but only $2.0/11.9 =
17$ per cent of the strontium.

\textbf{13.} $0.01\times40/1.1425 = 0.350$: 35 per cent.

\textbf{14.} From $R\,(40f + 0.75(1-f)) = 40fR_S + 0.75(1-f)R_D$:
$f = 0.75(R - R_D)/\bigl(40(R_S - R) + 0.75(R - R_D)\bigr)$, and the same with 200
and 10 for strontium.

\textbf{15.} $f = 0.75\times(-0.00029)/\bigl(40\times(-0.00137) + 0.75\times(-0.00029)\bigr)
= \num{-2.175e-4}/\num{-5.502e-2} = 0.0040$: 0.40 per cent.

\textbf{16.} $f = 10\times0.00128/\bigl(200\times0.01590 + 10\times0.00128\bigr) =
0.0128/3.1928 = 0.0040$: 0.40 per cent; the two agree.

\textbf{17.} Neodymium: $0.75\times(-0.00065)/\bigl(40\times(-0.00101) +
0.75\times(-0.00065)\bigr) = \num{-4.875e-4}/\num{-4.089e-2} = 0.0119$; strontium:
$10\times0.00336/\bigl(200\times0.01382 + 10\times0.00336\bigr) = 0.0336/2.7976 = 0.0120$.
Both give 1.2 per cent.

\textbf{18.} $f = 0.75\times(-0.00019)/\bigl(40\times(-0.00147) + 0.75\times(-0.00019)\bigr)
= 0.0024$, which predicts Sr \qty{10.46}{ppm} and $^{87}\mathrm{Sr}/^{86}\mathrm{Sr} =
0.70361$, against 0.70290 measured. Asha is not on the mantle--sediment mixing curve:
its source holds another component.

\textbf{19.} Old recycled oceanic crust, the HIMU type: its Nd isotopes need a Sm/Nd
somewhat below that of the \omterm{def:b2:isotope-geochemistry:reservoirs}{depleted mantle}, its Sr isotopes a low Rb/Sr, and its
lead (not measured here) a high U/Pb, all over one to two billion years.

\textbf{20.} $e^{\lambda_{147}\cdot2\times10^9} - 1 = 0.013133$. Sediment: $0.51141 -
0.110\times0.013133 = 0.509965$; CHUR: $0.512638 - 0.1966\times0.013133 = 0.510056$;
$\varepsilon_{\mathrm{Nd}}(2\,\mathrm{Ga}) = (0.509965/0.510056 - 1)\times10^4 = -1.8$.

\textbf{21.} $f = 0.110/0.1966 - 1 = -0.440$; since subduction $\varepsilon$ changed by
$25.0\times(-0.440)\times2.0 = -22.0$, so $\varepsilon(2\,\mathrm{Ga}) \approx -24.0 + 22.0
= -2.0$, within 0.2 unit of the exact value.

\textbf{22.} $e^{\lambda_{87}\cdot2\times10^9} - 1 = 0.028338$;
$0.72000 - 0.50\times0.028338 = 0.70583$.

\textbf{23.} At subduction the sediment had $\varepsilon_{\mathrm{Nd}} = -1.8$, close to the
$-4.7$ of a crust only 0.7 billion years old at the time, and well above the
strongly negative values of old crust: it fits the erosion of young crust, perhaps
mixed with some mantle-derived material. Its $^{87}\mathrm{Sr}/^{86}\mathrm{Sr}$ of
0.7058 at that time is also that of a young crustal rock.

\textbf{24.} The source of Vel, the most enriched island, holds
\textbf{\qty[text-series-to-math=true]{1.2}{\percent}} of recycled sediment.
\end{solution}
